% ============================================================
\section{Benchmark Evaluation and Optimality Gap}
\label{sec:experiments}
% ============================================================

This section reports the baseline benchmark study conducted to establish
a performance floor, motivate the need for learned allocation policies,
and quantify the improvement target available to a future reinforcement
learning agent. We run three deterministic greedy heuristics and a
manually computed full-information optimum across four benchmark
portfolio configurations of increasing complexity. All episodes are
executed through the Gymnasium-compatible environment described in
Section~\ref{sec:problem}, with the SQLite audit ledger providing a
complete per-period record of allocations, cash-flow events, and breach
flags for every episode.

% ------------------------------------------------------------
\subsection{Benchmark Configurations}
\label{subsec:configs}
% ------------------------------------------------------------

Four configurations are designed to test distinct structural aspects of
the environment. All configurations share the same contract parameters:
profit margin $\pi_i = 0.20$, S-curve shape $(a_i, b_i) = (1.5, 3.0)$,
retention rate $\nu_i = 0.05$, progress delay cap
$\delta^{\max}_i = 0.15$, termination tolerance $\theta_i = 2$,
finish delay cap $\Delta^{\max}_i = 2$, annual interest rate $r = 0.12$,
discount factor $\gamma = 0.95$, and budget tightness $\kappa = 3.0$.
Efficiency sequences $(\eta_{i,t})$ are fixed as custom ordinal draws to
ensure reproducibility across solvers. Table~\ref{tab:configs} summarises
the distinguishing features of each configuration.

\begin{table}[htbp]
  \centering\small
  \renewcommand{\arraystretch}{1.3}
  \setlength{\tabcolsep}{5pt}
  \begin{tabular}{lccccl}
    \toprule
    \textbf{Config ID} & $N$ & $\mathrm{BAC}_i$ & $D_i$ & $M_i$ &
    \textbf{Primary mechanic tested} \\
    \midrule
    TC-10SF-R   & 1 & 100 & 1 & 1 &
      Single project, single milestone, no advance payment \\
    TC-1LMAIF-R & 1 & 100 & 6 & 3 &
      Delayed start ($s_i=2$), advance payment, high recovery ($\rho_i=0.5$) \\
    TC-10MAIF-R & 1 & 100 & 6 & 3 &
      Degrading efficiency sequence, advance payment, termination risk \\
    TC-3XMAIF-R & 3 & 100/200/300 & 6 & 3 &
      Multi-project, staggered starts ($s_i = 0, 6, 7$), interacting cash flows \\
    \bottomrule
  \end{tabular}
  \caption{Benchmark configuration summary. $N$: number of projects;
    $\mathrm{BAC}_i$: budget at completion; $D_i$: planned duration
    (periods); $M_i$: number of non-advance milestones.}
  \label{tab:configs}
\end{table}

\paragraph{TC-10SF-R.}
The simplest configuration: one project, one milestone, no advance
payment, efficiency sequence $[0.91, 1.17, 0.83]$. The episode horizon
spans three active periods. This configuration tests the baseline
cash-flow mechanics in isolation --- advance payment, retention, and
milestone sequencing are all absent, so the only decision variable is
how aggressively to allocate in each period.

\paragraph{TC-1LMAIF-R.}
A single project with a delayed start at $s_i = 2$, three milestones,
advance payment rate $\phi_i = 0.20$ with advance trigger at progress
threshold 0.80, and a high advance recovery rate $\rho_i = 0.50$.
Efficiency alternates between above and below unity
$[0.91, 1.17, 0.83, 0.91, 1.17, 0.83]$. This configuration tests the
advance payment and recovery mechanic: the large recovery deduction from
milestone payments creates a temporary negative liquidity effect that
the allocation policy must anticipate.

\paragraph{TC-10MAIF-R.}
A single project with three milestones, advance payment $\phi_i = 0.20$,
and a monotonically \emph{degrading} efficiency sequence
$[0.95, 0.91, 0.55, 0.40, 0.30, 0.20]$. The sharp late-stage efficiency
drop means that allocations in later periods produce progressively less
progress per unit of spend, creating a strong incentive to front-load
work and certify milestones early. This configuration tests the
termination mechanic: with degrading efficiency it becomes increasingly
difficult to maintain progress above the $\delta^{\max}_i = 0.15$
tolerance floor.

\paragraph{TC-3XMAIF-R.}
The most complex configuration: three simultaneous projects with
$\mathrm{BAC} = [100, 200, 300]$, staggered starts at $s_i = [0, 6, 7]$,
and heterogeneous per-project efficiency sequences. The episode horizon
extends to period 11. Budget scarcity is real: the initial portfolio
budget is $B_0 = \kappa \sum_i \mathrm{BAC}_i = 3.0 \times 600 = 1{,}800$.
However, projects 1 and 2 start simultaneously at periods 6 and 7,
creating a large concurrent outflow window where allocations to three
active projects must be coordinated against a shared reserve while
advance payment outflows for projects 1 and 2 are incurred in rapid
succession. This configuration tests the full multi-project coordination
problem.

% ------------------------------------------------------------
\subsection{Greedy Heuristics and Optimal Reference}
\label{subsec:solvers}
% ------------------------------------------------------------

Three deterministic greedy heuristics are evaluated as defined in
Section~\ref{subsubsec:l3_baselines}. All three read directly from the
catchup and reach-plan fields of the observation vector and require no
training.

\textbf{Reach-Plan} ($\pi^{\mathrm{rp}}$) closes the progress gap to
the S-curve plan entirely each period, allocating
$x^{\mathrm{rp}}_{i,t}$ to every active project regardless of budget
scarcity or milestone proximity.

\textbf{Catchup-T} ($\pi^{\mathrm{cu}}$) is the minimum-spend policy:
it allocates only enough to keep the progress delay within
$\delta^{\max}_i$, preserving budget for later periods or other
projects.

\textbf{Catchup-Next-T} ($\pi^{\mathrm{cu}+1}$) is the most
risk-averse policy: it pre-emptively allocates enough this period so
that the projected delay \emph{next} period also stays within
$\delta^{\max}_i$, at the cost of higher early expenditure.

The \textbf{full-information optimum} ($\pi^{*}$) is computed
manually with complete knowledge of all efficiency realisations
$(\eta_{i,t})$ before the episode begins. It represents the allocation
that a perfect oracle with zero uncertainty would choose, and serves as
an upper-bound reference for the achievable episode NPV. Implementing a
tractable algorithmic solver --- whether a Linear Program or a
Dynamic Program --- is left as future work; the manual computation
suffices to establish the performance ceiling for the benchmark
configurations.

% ------------------------------------------------------------
\subsection{Results}
\label{subsec:results}
% ------------------------------------------------------------

Figure~\ref{fig:gap_plot} reports cumulative episode NPV for each
method across the four configurations. Table~\ref{tab:gap_results}
presents the numerical values and the gap between each heuristic and
the optimum.

\begin{figure}[htbp]
  \centering
  \includegraphics[width=\textwidth]{gap_plot.png}
  \caption{Cumulative episode NPV (episode $G = \sum_t \gamma^t \Delta B_t$)
    for three deterministic greedy heuristics and the full-information
    optimum across four benchmark configurations ordered by increasing
    portfolio complexity. The widening gap as complexity grows confirms
    that rule-based allocation is insufficient at portfolio scale and
    establishes a well-defined improvement target for learned agents.}
  \label{fig:gap_plot}
\end{figure}

\begin{table}[htbp]
  \centering\small
  \renewcommand{\arraystretch}{1.3}
  \setlength{\tabcolsep}{6pt}
  \begin{tabular}{lrrrrc}
    \toprule
    \textbf{Config} &
    $\pi^{*}$ &
    $\pi^{\mathrm{rp}}$ &
    $\pi^{\mathrm{cu}}$ &
    $\pi^{\mathrm{cu}+1}$ &
    \textbf{Best-heuristic gap} \\
    \midrule
    TC-10SF-R   &  $9.00$ &  $5.60$ &  $3.40$ &  $-6.07$ & $3.40$ \\
    TC-1LMAIF-R &  $1.49$ & $-13.21$ & $-13.21$ & $-27.17$ & $14.70$ \\
    TC-10MAIF-R & $-21.66$ & $-35.06$ & $-47.77$ & $-27.22$ & $5.56$ \\
    TC-3XMAIF-R & $12.40$ & $-59.75$ & $-62.38$ & $-107.96$ & $72.15$ \\
    \bottomrule
  \end{tabular}
  \caption{Episode NPV for all methods and configurations. Best-heuristic
    gap is $G(\pi^{*}) - \max_{\pi \in \{\pi^{\mathrm{rp}},
    \pi^{\mathrm{cu}}, \pi^{\mathrm{cu}+1}\}} G(\pi)$.}
  \label{tab:gap_results}
\end{table}

\paragraph{TC-10SF-R.}
On the simplest configuration the greedy baselines perform reasonably
well. Reach-Plan achieves $5.60$ against the optimal $9.00$, a gap of
$3.40$. Catchup-T follows at $3.40$. The risk-averse Catchup-Next-T
over-spends early and finishes at $-6.07$, demonstrating that excessive
caution is costly even in the single-project case when there is only one
milestone to unlock.

\paragraph{TC-1LMAIF-R.}
With a delayed project start, advance payment mechanics, and alternating
efficiency, all three heuristics produce negative NPV while the optimum
remains marginally positive at $1.49$. The best heuristic gap grows to
$14.70$. The advance payment recovery mechanism creates a liquidity trough
after the advance is disbursed at period 0 --- the optimal allocation
exploits the known efficiency sequence to time milestone certifications
precisely, avoiding unnecessary outflows during the recovery window.
The greedy policies, lacking this foresight, over-allocate during the
low-return periods and fail to offset the recovery deductions.

\paragraph{TC-10MAIF-R.}
The degrading efficiency sequence makes the allocation timing problem
acute. The optimal solution achieves $-21.66$ --- itself a loss, because
the late-stage efficiency collapse makes full project recovery
impossible. However, all three heuristics perform materially worse:
Reach-Plan at $-35.06$, Catchup-T at $-47.77$, and Catchup-Next-T at
$-27.22$. The best heuristic gap is $5.56$. Notably, Catchup-Next-T
outperforms both other heuristics here: its risk-averse lookahead
prevents the over-spending that Reach-Plan incurs in late periods where
efficiency has collapsed.

\paragraph{TC-3XMAIF-R.}
The multi-project configuration produces the starkest result. The
optimal allocation achieves $+12.40$: by coordinating across three
projects with staggered starts and exploiting the known efficiency
sequences, the full-information oracle turns a complex portfolio into a
profitable episode. Every greedy heuristic fails to do so: Reach-Plan
finishes at $-59.75$, Catchup-T at $-62.38$, and Catchup-Next-T
collapses to $-107.96$. The best-heuristic gap is $72.15$ NPV units ---
more than $20\times$ the gap on the simplest configuration. The sign
reversal is the critical observation: the optimal policy generates
positive returns while the best greedy baseline generates a substantial
loss.

% ------------------------------------------------------------
\subsection{Discussion}
\label{subsec:discussion}
% ------------------------------------------------------------

Three findings emerge from the benchmark study.

\textbf{The gap is monotonically increasing with complexity.}
The best-heuristic gap grows from $3.40$ on the single-project,
single-milestone baseline to $72.15$ on the three-project configuration,
a $21\times$ increase. This is not a marginal degradation; it reflects
a structural failure of greedy allocation under multi-project
coordination requirements. As projects overlap in time, their cash
flows interact through the shared budget $B_t$: an allocation decision
for project $i$ at period $t$ directly constrains what is available for
projects $j \neq i$, and the timing of advance payments and milestone
certifications across projects creates non-trivial sequencing
dependencies that no single-project greedy rule can resolve.

\textbf{No single heuristic dominates.}
Across the four configurations, each heuristic is best on at least one
case. Reach-Plan is best on TC-10SF-R; Catchup-Next-T is best on
TC-10MAIF-R; Catchup-T is best on TC-3XMAIF-R. This lack of a
uniformly dominant rule confirms that the optimal policy is adaptive:
the right allocation depends on the efficiency trajectory, the milestone
payment structure, and the number and phasing of concurrent projects in
a way that no fixed greedy rule can capture.

\textbf{A well-defined improvement target exists for learned agents.}
The full-information optimum $G(\pi^{*})$ provides an upper bound on
what a causal policy can aspire to. On TC-3XMAIF-R the gap between
the best achievable ($12.40$) and the best greedy ($-59.75$) is
$72.15$ NPV units, corresponding to a swing from a loss-making to a
profitable portfolio outcome. This gap is the value of intelligent
sequential allocation --- and it is precisely what a policy trained on
the CMDP environment of Section~\ref{sec:problem} is positioned to
close. Training such a policy and evaluating it against a formal Linear
Programming bound constitute the natural next step, addressed in
subsequent work.